\makeatletter
%Declaração das variáveis
\newcommand{\nomeuniversidade}[1]{\gdef\@nomeuniversidade{#1}}
\newcommand{\nomeinstituto}[1]{\gdef\@nomeinstituto{#1}}
\newcommand{\nomeinstitutoparceiro}[1]{\gdef\@nomeinstitutoparceiro{#1}}
\newcommand{\nomeprograma}[1]{\gdef\@nomeprograma{#1}}
\newcommand{\siglaprograma}[1]{\gdef\@siglaprograma{#1}}
\newcommand{\subtitulo}[1]{\gdef\@subtitulo{#1}}
\newcommand{\cidade}[1]{\gdef\@cidade{#1}}
\newcommand{\orientador}[1]{\gdef\@orientador{#1}}
\newcommand{\coorientador}[1]{\gdef\@coorientador{#1}}
\newcommand{\orientadorestitulo}[1]{\gdef\@orientadorestitulo{#1}}
\newcommand{\textofolharosto}[1]{\gdef\@textofolharosto{#1}}
\newcommand{\siglauniversidade}[1]{\gdef\@siglauniversidade{#1}}


%Local onde devem ser alteradas as informações gerais do trabalho
\nomeuniversidade{Universidade Federal do Rio de Janeiro}
\nomeinstituto{Instituto de Computação}
\nomeinstitutoparceiro{Instituto Tércio Pacitti de Aplicações e Pesquisas Computacionais}
\nomeprograma{Programa de Pós-Graduação em Informática}
\siglaprograma{PPGI}
\author{Nome do(a) Autor(a) do Trabalho}
\title{Título do Trabalho}
%Caso o trabalho não tenha subtítulo, basta deixar o conteudo do comando subtitulo em branco.
% Ex:\subtitulo{}
\subtitulo{: subtítulo do trabalho}
\cidade{Rio de Janeiro}
\date{2026}

\orientador{Nome do(a) Orientador(a)}
\coorientador{Nome do(a) Co-orientador(a)}
\orientadorestitulo{D.SC}
\textofolharosto{Dissertação de Mestrado apresentada ao \@nomeprograma, \@nomeinstituto\ e  \@nomeinstitutoparceiro, \@nomeuniversidade, como requisito parcial à obtenção do título de Mestre em Informática.}
\siglauniversidade{UFRJ}


\makeatother